\documentclass[10pt,a4paper]{article}
\usepackage[margin=0.9in]{geometry}
\usepackage{booktabs}
\usepackage{longtable}
\usepackage{array}
\usepackage{xcolor}
\usepackage{hyperref}
\usepackage{enumitem}
\usepackage{parskip}
\usepackage{amsmath}
\usepackage{amssymb}
\usepackage[T1]{fontenc}
\hypersetup{colorlinks=true,linkcolor=blue,urlcolor=blue,citecolor=blue}

\newcommand{\vok}{{\color{teal}\textbf{VERIFIED}}}
\newcommand{\vpc}{{\color{orange}\textbf{PARTIALLY CONTRADICTED}}}
\newcommand{\vnt}{{\color{gray}\textbf{NOT TESTED (wet-lab)}}}
\newcommand{\vin}{{\color{red}\textbf{INCONCLUSIVE}}}

\title{Detailed Replication Report:\\
Kalidasan et al.\ (2018) --- Putative Iron Acquisition Systems in
\textit{Stenotrophomonas maltophilia}}
\author{Ollie (agent-driven replication, BV-BRC pipeline)\\
BVBRC-06 wave, backfilled to 8-artifact standard on 2026-07-05}
\date{Original replication: 2026-05-10 \quad Backfill: 2026-07-05}

\begin{document}
\maketitle

\begin{abstract}
\noindent Kalidasan \textit{et al.}\ (2018, \textit{Molecules} 23:2048,
\href{https://doi.org/10.3390/molecules23082048}{doi:10.3390/molecules23082048})
combined in-silico RAST annotation of four complete \textit{S.\ maltophilia}
genomes (K279a, R551-3, D457, JV3) with wet-lab PCR (109 isolates), NanoString
nCounter gene expression (6 isolates $\times$ 17 targets), CAS/Arnow's
siderophore assays, and iron-source growth curves. This replication reproduces
the \textbf{in-silico slice only} using BV-BRC/RASTtk (the direct successor to
the RAST server used in 2018). All four genomes, both flagged iron-acquisition
subsystems, and all 17 functional targets in K279a were recovered; one paper
claim (DyP-type peroxidase being K279a-exclusive) is partially contradicted at
the gene level in modern RASTtk (the gene is present in all four strains,
though the subsystem \textit{assignment} is missing from all four --- suggesting
an ontology reshuffle, not a biology difference). \textbf{Verdict: REPLICATED}
for the computationally accessible component (7 of 22 total claims;
6/7 verified, 1/7 partially contradicted). The wet-lab 15/22 claims are
out-of-scope by definition. The critique section below is deliberately harsh
about how strong that "REPLICATED" verdict actually is.
\end{abstract}

\section{Paper Summary}
\subsection{Aim}
Identify putative iron acquisition systems in \textit{S.\ maltophilia} and the
iron sources this MDR opportunistic pathogen uses during iron starvation.

\subsection{Design}
Five layered components:
\begin{enumerate}[leftmargin=1.4em,itemsep=1pt]
  \item \textbf{In-silico}: RAST server annotation of 4 complete genomes
        (K279a, R551-3, D457, JV3).
  \item \textbf{PCR distribution}: 17 iron-acquisition-gene targets amplified
        across 103 clinical + 5 environmental + 1 reference (ATCC13637)
        = 109 isolates.
  \item \textbf{Expression}: NanoString nCounter Elements on 3 clinical (SM72,
        SM77, SM79) + 3 environmental (LMG10871, LMG10879, LMG11104) isolates
        under iron-depleted (100\,$\mu$M 2,2$'$-dipyridyl BHI) vs
        iron-repleted (BHI-DIP + 100\,$\mu$M FeCl$_3$) conditions.
  \item \textbf{Siderophore assays}: CAS agar diffusion (CASAD), liquid CAS,
        Arnow's colorimetric assay.
  \item \textbf{Iron-source utilization}: growth kinetics for SM77 + LMG10879
        with transferrin, lactoferrin, hemoglobin, hemin.
\end{enumerate}

\subsection{Headline findings (as reported in the paper)}
\begin{itemize}[leftmargin=1.4em,itemsep=1pt]
  \item Two putative iron acquisition subsystems dominate: (a) iron-siderophore
        sensor \& receptor system, (b) heme/hemin uptake and utilization.
  \item DyP-type peroxidase + ferritin-like protein oligomers subsystem
        \textbf{only in K279a}.
  \item FUR present in all four in-silico strains.
  \item 17 functional targets identified in K279a (Table 2 of paper).
  \item Clinical isolates $\gg$ environmental isolates for both PCR positivity
        and inducibility.
  \item NanoString: 6 targets significantly up-regulated under iron depletion
        --- FeSR (6.15$\times$, $p{=}0.023$), HmuT (12.21$\times$, $p{=}0.005$),
        Hup (5.46$\times$, $p{=}0.014$), ETFb (2.28$\times$, $p{=}0.010$),
        TonB (2.03$\times$, $p{<}0.01$), Fur (3.30$\times$, $p{=}0.003$).
  \item Arnow's: catechol-type siderophore.
  \item Liquid CAS: CS17 clinical 30.8\% vs LMG10879 environmental 4\%
        (only CS17 result reached $p{<}0.05$).
  \item Growth substrates: transferrin $>$ hemoglobin $\gg$ hemin, lactoferrin;
        decline from 72\,h onwards; positive control \textit{N.\ meningitidis}.
\end{itemize}

\section{Claims Table}
\begin{longtable}{@{}p{0.4cm}p{8cm}p{2.4cm}p{2.4cm}p{1.9cm}@{}}
\toprule
\# & Claim (paper section) & Type & Testable\,here? & Verdict \\
\midrule\endhead
C1  & 2 iron-acquisition subsystems shared across 4 strains (§2.1) & in-silico & YES & \vok \\
C2  & Iron siderophore sensor \& receptor system in all 4 (§2.1) & in-silico & YES & \vok \\
C3  & Heme/hemin uptake system in all 4 (§2.1) & in-silico & YES & \vok \\
C4  & DyP + ferritin-like oligomer subsystem K279a-only (§2.1) & in-silico & YES & \vpc \\
C5  & FUR present across all 4 (§2.1) & in-silico & YES & \vok \\
C6  & 17 functional targets identified in K279a Table~2 (§2.1) & in-silico & YES & \vok \\
C7  & 4 genomes with GenBank accessions AE016879 / CP001111 / HE798556 / CP002986 (Table~1) & metadata & YES & \vok \\
C8  & 109 isolates screened; clinical 100\% for Hyp1/Hup/ETFb/TonB/DyP/FUR (§2.2) & wet-lab & NO & \vnt \\
C9  & Environmental: only 8/17 targets amplified (§2.2) & wet-lab & NO & \vnt \\
C10 & FeSR 6.15$\times$ up, $p{=}0.023$ (§2.3) & wet-lab & NO & \vnt \\
C11 & HmuT 12.21$\times$ up, $p{=}0.005$ (§2.3) & wet-lab & NO & \vnt \\
C12 & Hup 5.46$\times$ up, $p{=}0.014$ (§2.3) & wet-lab & NO & \vnt \\
C13 & ETFb 2.28$\times$ up, $p{=}0.010$ (§2.3) & wet-lab & NO & \vnt \\
C14 & TonB 2.03$\times$ up, $p{<}0.01$ (§2.3) & wet-lab & NO & \vnt \\
C15 & Fur 3.30$\times$ up, $p{=}0.003$ (§2.3) & wet-lab & NO & \vnt \\
C16 & FeSS $-1.36\times$ (down, ns) (§2.3) & wet-lab & NO & \vnt \\
C17 & Clinical siderophore 30.8\% vs environmental 4\% liquid CAS (§2.4) & wet-lab & NO & \vnt \\
C18 & Arnow's: catechol-type siderophore (§2.4) & wet-lab & NO & \vnt \\
C19 & Transferrin $=$ max growth ($p{<}0.001$) (§2.5) & wet-lab & NO & \vnt \\
C20 & Hemoglobin second highest ($p{<}0.001$) (§2.5) & wet-lab & NO & \vnt \\
C21 & Growth decline from 72\,h onwards (§2.5) & wet-lab & NO & \vnt \\
C22 & Iron-depleted OD $1.007{\pm}0.276$ vs repleted $1.329{\pm}0.485$ (§2.4) & wet-lab & NO & \vnt \\
\bottomrule
\end{longtable}

\noindent\textbf{Coverage:} 7/22 = 31.8\% of paper claims are computationally
testable at all. Of those 7, 6 verified, 1 partially contradicted.
100\% coverage on the in-silico slice; 0\% on wet-lab claims (out of scope).

\section{Method (Replication Pipeline)}
\begin{enumerate}[leftmargin=1.4em,itemsep=2pt]
  \item Retrieved all four genomes from BV-BRC by GenBank accession:
        \texttt{522373.48} (K279a, 4{,}851{,}126\,bp, 1 contig);
        \texttt{391008.21} (R551-3, 4{,}573{,}969\,bp);
        \texttt{1163399.19} (D457, 4{,}769{,}156\,bp);
        \texttt{868597.17} (JV3, 4{,}544{,}477\,bp). Stored in
        \texttt{data/genome\_ids.json}.
  \item Queried BV-BRC "Iron acquisition and metabolism" subsystem class per
        genome via the public REST API and cached the full role/feature table
        (\texttt{data/iron\_subsystems\_all.json}).
  \item For K279a: mapped each of the 17 paper locus tags (\texttt{SMLT\_RS*}
        format, NCBI RefSeq) to BV-BRC \texttt{Smlt*} feature IDs and stored
        alongside the current RASTtk functional role
        (\texttt{data/k279a\_target\_mapping.json}).
  \item Comparative presence: for each target ran (a) PLfam species-level
        protein-family lookup and (b) keyword search on functional-role text
        across the other three genomes. Results in
        \texttt{data/comparative\_targets\_v2.json}.
  \item Narrative comparison to paper Table~2 in \texttt{analysis/subsystem\_comparison.md}
        and \texttt{analysis/gene\_presence\_comparison.md}.
\end{enumerate}

\noindent\textbf{Tool versions:} BV-BRC public REST API (as of 2026-05-10;
schema unchanged during backfill 2026-07-05); RASTtk annotation as delivered
by BV-BRC (successor to classic RAST 2018); PLfam species-level families as
served by BV-BRC. No local RASTtk re-run --- this replication trusts the
BV-BRC public annotation pipeline as-is.

\section{Results vs Paper}
\subsection{Subsystems (C1--C3)}
\begin{center}
\begin{tabular}{@{}lcccc@{}}
\toprule
Subsystem & K279a & R551-3 & D457 & JV3 \\
\midrule
Iron siderophore sensor \& receptor system         & \checkmark(8) & \checkmark(8) & \checkmark(5) & \checkmark(8) \\
Heme/hemin uptake \& utilization (GramPositives)   & \checkmark(3) & \checkmark(3) & \checkmark(3) & \checkmark(3) \\
\bottomrule
\end{tabular}
\end{center}
Two subsystems recovered in every strain, matching C1--C3.
Note the gene-count asymmetry in D457 (5, not 8) for the siderophore-receptor
system --- the paper does not resolve this counts-per-genome data cleanly.

\subsection{17 K279a targets (C6)}
All 17 targets mapped from paper locus tag ($\text{SMLT\_RS*}$) $\rightarrow$
BV-BRC (\text{Smlt*}) with matching functional role; see
\texttt{analysis/gene\_presence\_comparison.md} for the full table. 17/17 verified.

\subsection{Comparative distribution across 4 strains}
Via PLfam + keyword search, 17/17 targets present in all 4 strains.
HmuT PLfam did not resolve in R551-3 and JV3, but the "Heme ABC transporter,
cell surface" annotation is present under a different family --- caveat noted
in \texttt{gene\_presence\_comparison.md}. \textbf{Nothing here is a fresh
biology finding beyond what the paper already implies.}

\subsection{DyP claim (C4) --- partial contradiction}
Paper: "Encapsulating protein DyP-type peroxidase and ferritin-like protein
oligomers were only detected in K279a." \\
BV-BRC/RASTtk today: the \emph{subsystem} "Encapsulating protein DyP-type
peroxidase" is not emitted for \emph{any} of the 4 strains; the \emph{gene}
(\texttt{PLfam PLF\_40323\_00040048}, "Predicted dye-decolorizing peroxidase,
encapsulated subgroup") is present in \emph{all} 4 strains.

Two possibilities that this replication cannot distinguish:
\begin{enumerate}[label=(\alph*),leftmargin=1.4em]
  \item The 2018 classic RAST subsystem catalogue emitted the DyP subsystem
        for K279a but not the other three because their genomes lacked some
        secondary component (e.g.\ the ferritin-like partner) --- a real
        biology signal that modern RASTtk no longer surfaces.
  \item The 2018 assignment was an ontology artifact; the encapsulin/DyP
        subsystem was reorganised between RAST versions and current RASTtk
        emits neither the subsystem name nor a K279a-vs-others distinction.
\end{enumerate}
Distinguishing (a) vs (b) requires the 2018 SEED subsystem JSON for these four
genome IDs, which is not queryable in 2026 (the classic RAST server was
retired). This is called out again in Open Question Q1.

\subsection{FUR (C5)}
\texttt{fur} annotated in all 4 strains (K279a $=$ Smlt1986 with two
annotations; R551-3 = 1; D457 = 2; JV3 = 1). Verified.

\subsection{Wet-lab claims (C8--C22)}
Not attempted --- this is agent replication with a genomics API, no cultures.

\section{Per-Claim What-Worked / What-Didn't}
\textbf{Worked cleanly.} C1, C2, C3, C5, C6, C7 --- BV-BRC returned exactly
what the paper reports at the subsystem and locus level. The replication is
in the "look up in a database and compare to a table" regime; nothing surprising.

\textbf{Worked with a caveat.} HmuT PLfam missing in R551-3 / JV3 --- fell
back to a keyword search on the functional-role field. This is fine, but if
we were pushing hard we should have confirmed by amino-acid identity
($\ge 60\%$ over $\ge 80\%$ coverage) that the two candidate hits are true
orthologs. The replication script did not do this.

\textbf{Didn't work.} C4 (DyP K279a-only). At the gene level modern RASTtk
contradicts the paper; at the subsystem level modern RASTtk cannot even
support the paper's original assignment because the subsystem name is gone
from the emitted set. Marked partially contradicted; root cause not resolved.

\textbf{Not attempted.} C8--C22 wet-lab. No path to test in silico. But the
replication also did not push on \textit{indirect} in-silico tests that could
have partially derisked those claims:
\begin{itemize}[leftmargin=1.4em,itemsep=1pt]
  \item Presence of the 8 environmental-positive targets (Hyp1, Hup, ETFb,
        TonB, DyP, Fur, FeSR, Rp2) vs the 9 environmental-negative ones in
        \emph{all publicly sequenced} \textit{S.\ maltophilia} genomes ---
        would tell us whether the paper's "clinical $\gg$ environmental" split
        is a real clade signal or a sample-of-6 artifact.
  \item Sanity-check the primer sequences from paper Table~3 against the four
        BV-BRC genomes for perfect matches / degeneracy / dimer formation.
        The paper's 30--50\% amplification rates for Htp, ExbB, HmuV, FeSS in
        clinical isolates could partly be primer issues rather than gene
        absence.
\end{itemize}
Neither of these was done in the original replication window. They are logged
under Open Questions Q2 and Q3.

\section{Critique of the Replication (honest)}
\begin{itemize}[leftmargin=1.4em,itemsep=3pt]
\item \textbf{Verdict inflation risk.} "REPLICATED" here means "6 of 7 in-silico
claims recovered from a database that is a direct descendant of the tool the
authors used." That is a shallow bar. The paper's substantive scientific claims
--- differential expression of specific transporters, catechol siderophore
identity, transferrin as preferred iron source, clinical vs environmental
functional divergence --- are all wet-lab and are marked NOT TESTED. Calling
the whole paper REPLICATED at 31.8\% claim coverage overstates confidence.
A more honest label would be \textbf{"IN-SILICO SLICE REPLICATED; PAPER AS A WHOLE
NOT REPLICATED"}.

\item \textbf{RASTtk is not classic RAST.} The paper used the 2018 classic RAST
server; this replication used BV-BRC's RASTtk pipeline. They share ontology
lineage but diverge in subsystem call sets, protein-family models (PLfam vs
FIGfams), and role naming. The DyP mismatch (C4) is the visible tip of that
iceberg; there may be silent role reassignments across the other 16 targets
that we did not audit.

\item \textbf{No independent annotation.} The replication trusts BV-BRC's
served RASTtk annotations without re-running RASTtk (or Prokka / bakta) on the
raw genomes ourselves. If BV-BRC has a stale annotation snapshot for any of
these 4 genomes (they are all $\sim 2003$--$2011$ deposits with periodic
re-annotation), the replication just inherits that. No annotation-date check
was captured in the evidence.

\item \textbf{HmuT ortholog call is soft.} Where PLfam did not resolve, the
fallback was a keyword search on the functional-role string. Two proteins with
the same functional label are not necessarily orthologs; a proper reciprocal
best hit (BLAST or DIAMOND) was not run. Low-cost fix, not done.

\item \textbf{Locus-tag mapping is table-lookup, not proof.} The
$\text{SMLT\_RS*} \rightarrow \text{Smlt*}$ mapping used the NCBI GFF; if any
of those pairings are wrong, the "correct annotation" verdict propagates the
error silently. No cross-check by amino-acid sequence identity was performed.

\item \textbf{No siderophore / encapsulin biosynthetic-cluster analysis.} The
paper reports catechol-type siderophore production but does not identify the
molecule (e.g.\ enterobactin, dihydroxybenzoate-serine, stenobactin, etc.).
The replication could have run antiSMASH / PRISM on the four genomes to
predict siderophore biosynthetic gene clusters and give the paper's phenotypic
result a genotype. This was not attempted --- a real missed lever.

\item \textbf{The population-genomic extension was flagged but not run.}
BV-BRC hosts $\gtrsim 200$ \textit{S.\ maltophilia} genomes in 2026. Testing
the 17-target set across that full set would either strengthen the paper
(species-level signature) or contradict it (clinical/environmental subclade
structure). It is one query script away and was not done.

\item \textbf{Nougat parse is a stub.} Item 3 of the standard is a placeholder
noting sha256 + DOI; a later corpus sweep is required. This is called out in
\texttt{extraction/nougat.mmd}.

\item \textbf{Marker parse is a pdftotext fallback.} \texttt{extraction/marker.md}
is a layout-preserving \texttt{pdftotext} dump, not a real Marker parse; table
extraction quality (Tables 1--3 in particular) is degraded compared to true
Marker. Adequate for text re-read; not adequate for table re-mining.
\end{itemize}

\section{Verdict}
\textbf{REPLICATED} for the in-silico slice (6/7 claims verified, 1/7 partially
contradicted at gene-vs-subsystem level, 100\% of computationally accessible
claims tested). \textbf{NOT REPLICATED} for the full paper: 15 of 22 claims are
wet-lab and were not attempted (and cannot be attempted from a database alone).

\textbf{Confidence:} HIGH on the four-genome, 17-target roster match.
MEDIUM on the DyP subsystem story pending 2018-RAST-snapshot recovery.
NIL on all wet-lab expression / phenotype claims.

\section{Open Questions}
Each is grounded both in something Kalidasan et al.\ leave unresolved and
in something this replication surfaced. All five are intended as
research prompts, not paper "future work" boilerplate.

\subsection*{Q1 --- DyP peroxidase: real K279a specificity or ontology artifact?}
\textbf{Basis.} This replication directly contradicts C4 at the gene level:
the DyP gene is present in all four strains under modern RASTtk, but the paper
reported an encapsulin/DyP subsystem exclusive to K279a. The paper itself
never performed a DyP activity assay to confirm K279a-specific function.

\textbf{Next steps.} (1) Re-run classic RAST offline via SEEDtk on the four
assemblies to test whether the 2018 subsystem call reproduces. (2) Search
each genome for the encapsulin shell protein (Pfam PF04454) and ferritin-like
partner (Pfam PF00210) independently of DyP --- the K279a-only signal, if
biological, may live in the shell/partner rather than DyP itself.
(3) Wet-lab: H$_2$O$_2$-dependent ABTS oxidation activity on iron-depleted
lysates of all four strains, gel-filtration into $>200$\,kDa fraction to test
for encapsulation. (4) Cryo-EM / negative-stain EM on K279a to visualise
encapsulin nanocompartments.

\subsection*{Q2 --- What actually is the S.\ maltophilia catechol siderophore?}
\textbf{Basis.} The paper reports catechol-class by Arnow's ("data not shown")
but never identifies the molecule; the paper's own literature citation to
"hydroxamate-type ornibactin" is inconsistent with its catechol result and the
inconsistency is unresolved. This replication also did not run any biosynthetic-
cluster prediction on the four genomes --- a real missed lever.

\textbf{Next steps.} (1) antiSMASH v7/8 across all 4 genomes; enumerate
catecholate NRPS clusters and check for entA/B/E/F homologs. (2) LC-MS/MS on
iron-depleted CS17 supernatant (per paper conditions) after XAD-16 solid-phase
extraction; compare vs enterobactin, 2,3-DHB, and DHB-serine standards.
(3) Delete the top-scoring NRPS cluster in K279a and repeat CAS/Arnow's
--- loss of activity confirms the source cluster. (4) Cross-strain: does
R551-3 have a truncated/absent version of the same cluster, consistent with
its weaker CAS signal?

\subsection*{Q3 --- Clinical $\gg$ environmental: real ecotype, sample-size, or primer artifact?}
\textbf{Basis.} PCR-negative rates of 26.9\%--53.8\% in clinical isolates for
HmuV/FeSS/FCR/Htp, and 0\% amplification of 9-of-17 targets in only 5
environmental isolates, are equally consistent with (a) real gene loss in an
environmental subclade, (b) sample-size noise (n$=$5), or (c) primer failure
(paper Table~3 primers were designed on the 4 in-silico genomes; SNP density
between clinical and environmental \textit{S.\ maltophilia} clades is high).
The paper does not attempt an in-silico primer coverage check; nor did this
replication.

\textbf{Next steps.} (1) Pull all $\gtrsim 200$ public \textit{S.\ maltophilia}
BV-BRC genomes; test the 17-target set for statistically significant
clinical/environmental gene-presence differences (Fisher's exact + BH).
(2) In-silico primer coverage: align each Table~3 primer pair against the
full genome panel; count perfect matches, 3$'$-end mismatches, and predicted
amplicon size. (3) Redesign degenerate primers where coverage $<95\%$ and
reanalyze the 109-isolate panel. (4) Core-genome phylogeny (ParSNP/Roary) to
test whether Kalidasan's environmental isolates cluster in a known Sgn3/Sgn4
subclade with characterised gene losses.

\subsection*{Q4 --- Transferrin as top iron source without transferrin receptors: mechanism?}
\textbf{Basis.} The paper's Fig~3 shows transferrin gives maximal growth
($p<0.001$), yet the paper's own 17-target Table~2 (which this replication
re-confirmed) contains \emph{no} TbpA/TbpB/LbpA/LbpB homologs. The mechanism
by which \textit{S.\ maltophilia} liberates iron from Tf is unaddressed
in the paper's Discussion, which waves at "siderophore could potentially
scavenge iron-saturated Tf" but does not test it. Positive control
\textit{N.\ meningitidis} \emph{does} have TbpAB/LbpAB and is not comparable.

\textbf{Next steps.} (1) HMMER search of all 4 genomes for TbpA (PF01298),
TbpB (PF07298), LbpA/LbpB, and any TonB-dependent receptor with predicted
Tf-binding specificity. (2) Growth on holo-Tf $\pm$ apoferritin controls;
if only iron-saturated citrate supports growth, \textit{S.\ maltophilia} is
not using Tf directly. (3) Knock out the secreted metalloproteases (StmPr1/2/3)
in K279a and repeat Tf-growth assay --- proteolytic iron release hypothesis.
(4) Cross-reference Garcia et al.\ 2015 (\textit{Res Microbiol}) on
\textit{S.\ maltophilia} lack of Tf receptors.

\subsection*{Q5 --- Are any of the 6 iron-induced targets required for virulence?}
\textbf{Basis.} The paper's own concluding sentence: "\textit{These data need
to be further confirmed through several knockout studies}" --- the authors
themselves flag it. As of 2026, no follow-up study has closed the loop for
the 6-target set (FeSR, HmuT, Hup, ETFb, TonB, Fur). This replication cannot
close it in silico; but it does re-confirm the genes exist in K279a and are
ready for mutagenesis. Also: the ETFb induction (2.28$\times$) is unexplained
mechanistically --- ETFb is an electron-transfer flavoprotein, not an
obviously iron-acquisition gene.

\textbf{Next steps.} (1) Markerless in-frame deletion mutants in K279a for
each target (pEX18Tc adaptation, sacB counter-selection). (2) Phenotype panel:
BHI-DIP growth, CAS/Arnow's siderophore production, biofilm formation, serum
survival (50\% human serum), \textit{Galleria mellonella} or mouse-lung
colonisation. (3) In-trans complementation to confirm phenotype specificity.
(4) For ETFb specifically: dual transcriptomics + metabolomics under iron
depletion in wildtype vs $\Delta$etfB --- test electron-transfer-vs-iron
hypothesis.


\end{document}
